\documentclass[12pt]{article}

\usepackage{sbc-template}
\usepackage{graphicx,url}
\usepackage[brazil]{babel}
\usepackage[utf8]{inputenc}
\usepackage[T1]{fontenc}
\usepackage{amsmath}
\usepackage{xcolor}

\sloppy

% ---------------------------------------------------------------------------
% Números que dependem dos experimentos. Enquanto codigo/analisar.py não
% rodar, aparecem como [a medir]. Depois, resultados/macros.tex os substitui.
% ---------------------------------------------------------------------------
\newcommand{\pendente}{\textbf{\textcolor{red}{[a medir]}}}
\newcommand{\breakeven}{\pendente}
\newcommand{\breakevenEn}{\pendente}
\newcommand{\nEstimado}{\pendente}
\newcommand{\banda}{\pendente}
\newcommand{\speedupEteMax}{\pendente}
\newcommand{\speedupEteMin}{\pendente}
\newcommand{\speedupKernelMax}{\pendente}
\newcommand{\nMax}{4096}
\newcommand{\nMin}{64}
\newcommand{\fracTransfMin}{\pendente}
\newcommand{\fracTransfMax}{\pendente}
\newcommand{\cvMax}{\pendente}
\newcommand{\gflopsCpu}{\pendente}
\newcommand{\gflopsGpu}{\pendente}
\newcommand{\nSeq}{1024}
\newcommand{\speedupSeqKernel}{\pendente}
\newcommand{\speedupSeqEte}{\pendente}
\newcommand{\speedupRefEteSeqN}{\pendente}
\newcommand{\ganhoPinned}{\pendente}
\InputIfFileExists{resultados/macros.tex}{}{}

% Trechos de interpretação que precisam ser conferidos com os dados reais.
% Para a versão de submissão, troque \revisartrue por \revisarfalse.
\newif\ifrevisar \revisarfalse
\newcommand{\revisar}[1]{\ifrevisar{\color{orange!70!black}#1}\else#1\fi}

\newcommand{\figura}[2]{\IfFileExists{figuras/#1.pdf}%
  {\includegraphics[width=#2]{figuras/#1.pdf}}%
  {\fbox{\parbox[c][3.5cm][c]{\dimexpr#2-2\fboxsep-2\fboxrule}{\centering\small Figura gerada por\\\texttt{codigo/analisar.py}}}}}

\title{Quando a GPU compensa? Análise experimental de desempenho\\
entre CPU e GPU na multiplicação de matrizes}

% Se a revisão for double-blind, remova nomes, instituição e e-mails.
\author{Luís Filipe Alves Silva Santos\inst{1}, Matheus Costa da Rocha\inst{1},\\
Ronierison de Souza Maciel\inst{1}}

\address{Bacharelado em Ciência da Computação -- CESAR School\\
  Recife -- PE -- Brasil
  \email{\{lfass,mcr\}@cesar.school, rsm4@cesar.school}
}

\begin{document}

\maketitle

\begin{abstract}
GPU speedups over CPUs are often reported as large factors, but such figures
frequently compare an optimized GPU kernel against unoptimized single-threaded
CPU code and ignore host--device data transfers. This paper experimentally
investigates when a GPU actually pays off for single-precision dense matrix
multiplication. We compare naive and optimized multithreaded CPU
implementations (OpenMP and OpenBLAS) with three CUDA implementations (naive,
shared-memory tiling and cuBLAS), measuring both kernel-only and end-to-end
times, with repeated runs and 95\% confidence intervals. The GPU only
outperforms the best CPU implementation, end to end, beyond a break-even point
\breakevenEn{}, reaching a speedup of \speedupEteMax$\times$ at
$N=\nMax$, far below the \speedupSeqKernel$\times$ obtained when the kernel is
compared with naive sequential code.
\end{abstract}

\begin{resumo}
Ganhos de desempenho da GPU sobre a CPU costumam ser divulgados como fatores
expressivos, mas muitas vezes comparam um kernel de GPU otimizado com código
sequencial não otimizado na CPU e ignoram a transferência de dados entre as
memórias. Este trabalho investiga experimentalmente quando a GPU de fato
compensa na multiplicação de matrizes densas em precisão simples. Comparamos
implementações de CPU ingênua e otimizadas multithread (OpenMP e OpenBLAS) com
três implementações em CUDA (ingênua, blocada em memória compartilhada e
cuBLAS), medindo os tempos só do kernel e ponta a ponta, com repetições e
intervalos de confiança de 95\%. A GPU só supera a melhor implementação de
CPU, ponta a ponta, a partir de uma ordem de matriz \breakeven, chegando a um
speedup de \speedupEteMax$\times$ em $N=\nMax$, bem abaixo dos
\speedupSeqKernel$\times$ obtidos ao comparar o kernel com o código sequencial
ingênuo.
\end{resumo}

% ===========================================================================
\section{Introdução}
% ===========================================================================

Unidades de processamento gráfico (GPUs) deixaram de ser dispositivos
dedicados apenas à renderização e se tornaram processadores de propósito geral
voltados para vazão, com milhares de unidades aritméticas e memória de alta
largura de banda~\cite{owens2008}. Modelos de programação como o
CUDA~\cite{nickolls2008} tornaram esse hardware acessível, e a multiplicação
de matrizes densas, com $2N^3$ operações de ponto flutuante sobre $3N^2$
elementos, virou o exemplo clássico de carga que se beneficia dele.

É comum encontrar na literatura e na divulgação técnica ganhos de 10 a
100~vezes da GPU sobre a CPU. Esses números, porém, dependem fortemente de
\emph{como} a comparação é feita. Lee et al.~\cite{lee2010} mostraram
que, quando o código de CPU recebe o mesmo cuidado de otimização que o de GPU
(uso de todos os núcleos, vetorização SIMD e blocagem em cache), a diferença
média cai para poucas vezes. Gregg e Hazelwood~\cite{gregg2011}
argumentaram que não se pode comparar CPU e GPU sem dizer onde os dados estão:
se as matrizes nascem e terminam na memória principal, o tempo de cópia pelo
barramento PCIe faz parte do custo e pode dominar o tempo total.

Este trabalho parte dessas duas críticas e formula uma pergunta prática:
\textbf{a partir de que tamanho de problema a GPU passa a compensar, de ponta a
ponta, em relação a uma CPU multithread otimizada?} Para respondê-la, a
multiplicação de matrizes quadradas em precisão simples é executada em seis
implementações (três de CPU e três de GPU), em tamanhos de $N=\nMin$ a
$N=\nMax$, separando o tempo do kernel do tempo de transferência e tratando
os resultados com repetições e intervalos de confiança.

As contribuições são: (i) uma metodologia de medição justa e reprodutível,
com código aberto, que inclui a transferência CPU--GPU e usa código
otimizado também na CPU; (ii) a identificação experimental do ponto de
equilíbrio em que a GPU passa a vencer, comparado com um modelo analítico
simples; e (iii) uma quantificação de quanto o speedup reportado muda conforme
a linha de base escolhida.

O restante do texto está organizado assim: a Seção~\ref{sec:relacionados}
discute os trabalhos relacionados; a Seção~\ref{sec:metodologia} descreve
implementações, ambiente e protocolo de medição; a
Seção~\ref{sec:resultados} apresenta e discute os resultados; a
Seção~\ref{sec:ameacas} trata das ameaças à validade; e a
Seção~\ref{sec:conclusao} conclui.

% ===========================================================================
\section{Fundamentação e Trabalhos Relacionados}\label{sec:relacionados}
% ===========================================================================

\textbf{GPUs e CUDA.} Owens et al.~\cite{owens2008} descrevem a
evolução das GPUs para computação de propósito geral e as características
que favorecem aplicações com paralelismo de dados abundante, como a
multiplicação de matrizes. Nickolls et al.~\cite{nickolls2008}
apresentam o modelo CUDA, em que o programador organiza a computação em uma
grade de blocos de threads que cooperam por memória compartilhada e barreiras
de sincronização. Hwu et al.~\cite{hwu2022} detalham a técnica de
blocagem (\emph{tiling}) em memória compartilhada, que reduz o tráfego com a
memória global ao reaproveitar cada elemento lido por vários threads; essa é a
base da nossa implementação intermediária de GPU.

\textbf{Multiplicação de matrizes otimizada.} Em CPU, implementações de alto
desempenho seguem a organização proposta por
Goto e van de Geijn~\cite{goto2008}, que particiona as matrizes em blocos
ajustados a cada nível da hierarquia de cache e usa microkernels vetorizados;
essa organização deu origem à GotoBLAS e, por extensão, à OpenBLAS usada aqui.
Em GPU, Volkov e Demmel~\cite{volkov2008} mostraram, por
microbenchmarks, como ajustar a multiplicação de matrizes para aproximá-la do
pico do hardware, resultados incorporados a bibliotecas como a cuBLAS. O
modelo Roofline~\cite{williams2009} relaciona o desempenho atingível à
intensidade aritmética (operações por byte movido) e é útil para entender por
que a multiplicação de matrizes, cuja intensidade cresce com $N$, tende a ser
limitada por computação em matrizes grandes.

\textbf{Comparações justas entre CPU e GPU.} O trabalho mais próximo do nosso é
o de Lee et al.~\cite{lee2010}, que avaliou 14 kernels de vazão em
uma CPU e uma GPU da época, otimizando ambos, e concluiu que a vantagem média
da GPU era muito menor que os 100$\times$ frequentemente divulgados.
Gregg e Hazelwood~\cite{gregg2011} complementam esse argumento ao
mostrar que o tempo de transferência entre memórias muda a conclusão sobre qual
processador é melhor e propõem que toda comparação informe onde os dados
residem. Nosso trabalho aplica as duas recomendações a um único kernel, em
hardware atual e acessível a um estudante, e vai além ao localizar
experimentalmente o tamanho de problema a partir do qual a GPU passa a vencer
de ponta a ponta.

% ===========================================================================
\section{Metodologia}\label{sec:metodologia}
% ===========================================================================

\subsection{Pergunta de pesquisa e hipóteses}

A pergunta central é: \emph{para multiplicação de matrizes densas quadradas em
FP32 com os dados inicialmente na memória da CPU, a partir de que ordem $N$ a
GPU supera, de ponta a ponta, a melhor implementação multithread de CPU?}
Formulamos três hipóteses:
\textbf{H1:} para $N$ pequeno, a CPU otimizada vence, porque os custos fixos de
transferência e de lançamento de kernel na GPU superam o tempo de computação;
\textbf{H2:} existe um ponto de equilíbrio $N^*$ a partir do qual a GPU vence,
e o speedup ponta a ponta cresce com $N$, já que a computação cresce com $N^3$
e a transferência com $N^2$;
\textbf{H3:} o speedup medido contra uma CPU sequencial ingênua, considerando
só o kernel, é ordens de grandeza maior que o medido contra a CPU otimizada
de ponta a ponta.

\textbf{Modelo analítico.} Sejam $P_c$ e $P_g$ as vazões (FLOP/s) da CPU e do
kernel de GPU, e $B$ a largura de banda efetiva do barramento. Copiar $A$ e
$B$ para a GPU e trazer $C$ de volta move $12N^2$ bytes em FP32. Desprezando
latências fixas, a GPU vence quando
\begin{equation}
\frac{2N^3}{P_c} > \frac{12N^2}{B} + \frac{2N^3}{P_g}
\quad\Longrightarrow\quad
N > N^* = \frac{6\,P_c}{B\,(1 - P_c/P_g)}.
\label{eq:nstar}
\end{equation}
A Equação~\ref{eq:nstar} mostra que o ponto de equilíbrio é proporcional à
vazão da CPU: quanto melhor o código de CPU, maior o problema necessário para
a GPU compensar. Na Seção~\ref{sec:resultados}, $N^*$ é estimado com $P_c$,
$P_g$ e $B$ medidos e comparado com o valor observado.

\subsection{Implementações}

Todas as versões calculam $C = AB$ com matrizes quadradas em FP32, armazenadas
por linhas, preenchidas com valores pseudoaleatórios uniformes em $[-1,1]$ a
partir da mesma semente, de modo que CPU e GPU recebem exatamente as mesmas
entradas.

\textbf{CPU ingênua (1 thread):} três laços aninhados na ordem $ijk$, sem
paralelismo; representa a linha de base usada em muitas comparações e só é
medida até $N=\nSeq$, pois fica lenta demais acima disso.
\textbf{CPU OpenMP blocada:} blocos de $64\times64$, ordem $ikj$ no laço
interno (acesso contíguo e vetorizável com \texttt{omp simd}) e paralelismo
sobre faixas de linhas de $C$ com OpenMP, sem escritas compartilhadas entre
threads.
\textbf{CPU OpenBLAS:} chamada \texttt{cblas\_sgemm} da OpenBLAS, multithread,
como referência do estado da arte em CPU~\cite{goto2008}.
\textbf{GPU ingênua:} um thread por elemento de $C$, lendo $A$ e $B$
diretamente da memória global, em blocos de $32\times32$ threads.
\textbf{GPU blocada:} blocagem em memória compartilhada com
\emph{tiles} de $32\times32$~\cite{hwu2022}.
\textbf{GPU cuBLAS:} \texttt{cublasSgemm} com o modo de precisão padrão
(sem TF32), para que a aritmética seja comparável à da CPU.

A corretude de cada execução é verificada recalculando 64 elementos de $C$,
escolhidos aleatoriamente, em precisão dupla; o erro relativo, normalizado por
$\sum_k|a_{ik}b_{kj}|$, deve ficar abaixo de $10^{-3}$.

\subsection{Ambiente experimental}

A Tabela~\ref{tab:ambiente} resume o hardware e o software. Os dados são
coletados automaticamente pelo script de experimentos.

\begin{table}[ht]
\centering
\caption{Ambiente experimental.}
\label{tab:ambiente}
\small
\begin{tabular}{ll}
\hline
CPU & Intel Core i7-9700F, 8 núcleos/8 threads, 3,0~GHz (turbo até 4,7~GHz), L3 de 12~MiB \\
Memória principal & 16~GB DDR4 (7,7~GiB disponíveis no WSL2) \\
GPU & NVIDIA GeForce RTX 3060, 12~GB GDDR6, 28 SMs \\
Barramento & PCIe 3.0 x16 \\
Sistema operacional & Ubuntu 26.04.1 LTS sob WSL2 (kernel 6.18), Windows 11 Pro \\
Driver & NVIDIA 621.24 \\
Compiladores & GCC 15.2 (\texttt{-O3 -march=native -fopenmp}); CUDA 12.4 (\texttt{-O3}) \\
Bibliotecas & OpenBLAS 0.3.32 (pthreads); cuBLAS 12.4 \\
\hline
\end{tabular}
\end{table}

\subsection{Métricas e protocolo de medição}

Para cada implementação e cada $N \in \{64, 128, \dots, \nMax\}$ (potências de
2), executamos 5 repetições de aquecimento, descartadas, seguidas de 30
repetições medidas, em processos separados por configuração. O número de
threads de CPU é igual ao de núcleos físicos. Na versão OpenMP, as threads são
fixadas aos núcleos (\texttt{OMP\_PROC\_BIND=close},
\texttt{OMP\_PLACES=cores}). Na OpenBLAS, essa fixação é desativada
(\texttt{OMP\_PROC\_BIND=false}): com ela ativa, o \emph{runtime} do OpenMP
prende a thread principal a um único núcleo e as threads criadas pela OpenBLAS
herdam essa restrição, de modo que a biblioteca roda, na prática, em um só
núcleo. Em testes preliminares, isso a deixou cerca de 6 vezes mais lenta
(aproximadamente 120 contra 760~GFLOP/s em $N=2048$), um exemplo de como um
detalhe de configuração pode subestimar a linha de base de CPU.

Na CPU, o tempo é medido com \texttt{clock\_gettime(CLOCK\_MONOTONIC)} em
torno da chamada de multiplicação. Na GPU, cada repetição registra quatro
tempos: cópia de $A$ e $B$ para a GPU, kernel e cópia de $C$ de volta, com
eventos CUDA, e o tempo \textbf{ponta a ponta}, com relógio de parede, que
inclui as três etapas e a sincronização final. A criação do contexto CUDA, a
alocação na GPU e a criação do \emph{handle} da cuBLAS ficam fora da medição,
por serem custos únicos amortizados numa aplicação real; a memória da CPU é
paginável (\texttt{malloc}), o caso mais comum. Como experimento adicional,
repetimos a cuBLAS com memória fixada (\emph{pinned}).

Reportamos, para cada configuração, a \textbf{mediana} das repetições, com
intervalo de confiança de 95\% obtido por \emph{bootstrap} (2.000
reamostragens), além do coeficiente de variação. Optamos pela mediana porque
observamos picos isolados de até 25 vezes o tempo típico, causados por
interferência do sistema operacional hospedeiro: como a OpenBLAS ocupa todos os
núcleos, basta um deles ser interrompido para que toda a multiplicação espere.
A média seria dominada por esses picos, enquanto a mediana representa o
desempenho típico~\cite{jain1991}. O desempenho é expresso em GFLOP/s
$= 2N^3/(t \cdot 10^9)$ e o speedup como a razão entre as medianas da melhor
CPU e da GPU, com intervalo de confiança obtido por \emph{bootstrap} com 4.000
reamostragens.

% ===========================================================================
\section{Resultados e Discussão}\label{sec:resultados}
% ===========================================================================

\revisar{Todas as medições passaram na verificação de corretude. A
variabilidade foi alta nos tamanhos pequenos, em que cada execução dura poucos
microssegundos e qualquer interrupção do sistema hospedeiro pesa (coeficiente
de variação de até \cvMax\%), e menor nos tamanhos maiores; por isso, os
resultados usam a mediana, como descrito na Seção~\ref{sec:metodologia}.} A
Tabela~\ref{tab:resultados} traz os tempos medianos das principais
configurações e as Figuras~\ref{fig:tempo-speedup} e~\ref{fig:decomp}
resumem o comportamento.

\begin{table}[ht]
\centering
\caption{Tempo mediano por multiplicação (mediana $\pm$ meia-largura do IC
95\% por \emph{bootstrap}, em \% da mediana). ``p.~a~p.'' = ponta a ponta, incluindo transferências.}
\label{tab:resultados}
\IfFileExists{resultados/tabela.tex}%
  {\resizebox{\textwidth}{!}{\begin{tabular}{rrrrrrr}
\hline
$N$ & CPU 1T & CPU OMP & OpenBLAS & GPU tiled (p.~a~p.) & cuBLAS (kernel) & cuBLAS (p.~a~p.) \\ \hline
64 & 148.0\,$\mu$s $\pm$0.1\% & 19.3\,$\mu$s $\pm$0.3\% & 5.9\,$\mu$s $\pm$0.0\% & 231.9\,$\mu$s $\pm$6.2\% & 8.2\,$\mu$s $\pm$3.4\% & 186.9\,$\mu$s $\pm$4.8\% \\
128 & 1.62\,ms $\pm$3.7\% & 75.4\,$\mu$s $\pm$0.1\% & 10.0\,$\mu$s $\pm$1.5\% & 187.6\,$\mu$s $\pm$1.2\% & 7.2\,$\mu$s $\pm$7.1\% & 190.5\,$\mu$s $\pm$5.1\% \\
256 & 14.37\,ms $\pm$0.4\% & 346.6\,$\mu$s $\pm$3.0\% & 47.4\,$\mu$s $\pm$0.4\% & 329.3\,$\mu$s $\pm$0.6\% & 42.3\,$\mu$s $\pm$1.3\% & 303.2\,$\mu$s $\pm$1.1\% \\
512 & 199.71\,ms $\pm$6.3\% & 4.15\,ms $\pm$8.4\% & 322.9\,$\mu$s $\pm$449.0\% & 937.4\,$\mu$s $\pm$1.1\% & 78.3\,$\mu$s $\pm$0.4\% & 648.1\,$\mu$s $\pm$1.8\% \\
1024 & 1.40\,s $\pm$0.6\% & 19.65\,ms $\pm$2.7\% & 2.49\,ms $\pm$8.4\% & 4.21\,ms $\pm$1.2\% & 306.5\,$\mu$s $\pm$3.3\% & 2.00\,ms $\pm$5.4\% \\
2048 & -- & 211.70\,ms $\pm$1.7\% & 22.47\,ms $\pm$1.9\% & 25.08\,ms $\pm$2.5\% & 2.18\,ms $\pm$0.2\% & 8.31\,ms $\pm$0.8\% \\
4096 & -- & 2.17\,s $\pm$0.6\% & 223.48\,ms $\pm$3.4\% & 166.65\,ms $\pm$0.2\% & 17.05\,ms $\pm$0.3\% & 42.45\,ms $\pm$4.8\% \\
\hline
\end{tabular}
}}%
  {\fbox{\parbox{0.8\textwidth}{\centering\small Tabela gerada por
  \texttt{codigo/analisar.py}}}}
\end{table}

\begin{figure}[ht]
\centering
\begin{minipage}[t]{0.49\textwidth}\centering
\figura{tempo}{\linewidth}\\
\small (a) Tempo por multiplicação
\end{minipage}\hfill
\begin{minipage}[t]{0.49\textwidth}\centering
\figura{speedup}{\linewidth}\\
\small (b) Speedup da cuBLAS sobre a OpenBLAS
\end{minipage}
\caption{(a) Tempo mediano em função de $N$; as curvas de GPU em linha
contínua são ponta a ponta e a tracejada é só o kernel da cuBLAS. (b) Speedup
da GPU sobre a melhor CPU; valores abaixo de 1 indicam que a CPU é mais rápida.
Barras de erro: IC 95\%.}
\label{fig:tempo-speedup}
\end{figure}

\textbf{Pequenos problemas (H1).} \revisar{Para $N=\nMin$, a GPU é
\emph{mais lenta} que a OpenBLAS de ponta a ponta (speedup de
\speedupEteMin$\times$). Nesse regime, como mostra a
Figura~\ref{fig:decomp}, cerca de \fracTransfMin\% do tempo da GPU é gasto
copiando dados, e o próprio kernel dura poucos microssegundos, próximo do
custo fixo de lançamento. A hipótese H1 é, portanto, confirmada.}

\textbf{Ponto de equilíbrio (H2).} \revisar{A GPU passa a vencer a melhor CPU
a partir de um ponto de equilíbrio situado \breakeven, e o speedup ponta a ponta cresce com $N$,
chegando a \speedupEteMax$\times$ em $N=\nMax$. O modelo da
Equação~\ref{eq:nstar}, alimentado com a vazão medida da CPU
(\gflopsCpu~GFLOP/s), do kernel da cuBLAS (\gflopsGpu~GFLOP/s) e com a banda
efetiva medida do barramento (\banda~GB/s), prevê $N^* \approx \nEstimado$,
valor que cai dentro do intervalo observado, o que indica que o modelo,
apesar de simples, captura bem o compromisso entre computação e transferência. A fração do tempo gasta em transferência cai para
\fracTransfMax\% em $N=\nMax$, coerente com a razão $N^3/N^2$, mas continua
relevante: o speedup só do kernel (até \speedupKernelMax$\times$) é bem maior
que o ponta a ponta, o que reforça o argumento de
Gregg e Hazelwood~\cite{gregg2011}. Com memória \emph{pinned}, o tempo
ponta a ponta da cuBLAS em $N=\nMax$ cai \ganhoPinned\%.}

\begin{figure}[ht]
\centering
\figura{decomposicao}{0.6\textwidth}
\caption{Decomposição do tempo ponta a ponta da cuBLAS em cópia para a GPU,
kernel e cópia de volta, em função de $N$.}
\label{fig:decomp}
\end{figure}

\textbf{A escolha da linha de base (H3).} \revisar{Em $N=\nSeq$, comparar o
kernel da cuBLAS com a CPU sequencial ingênua resulta em um speedup de
\speedupSeqKernel$\times$; incluindo as transferências, o valor cai para
\speedupSeqEte$\times$; e contra a OpenBLAS multithread, ponta a ponta, é de
apenas \speedupRefEteSeqN$\times$. O mesmo experimento pode, portanto,
sustentar números que diferem em ordens de grandeza, dependendo apenas da
linha de base e do que se inclui na medição, em linha com
Lee et al.~\cite{lee2010}. A hipótese H3 é confirmada.}

\textbf{Implementações próprias de GPU.} \revisar{A versão blocada em memória
compartilhada é mais rápida que a ingênua em todos os tamanhos, mas o ganho é
modesto (cerca de 1,2 a 1,3 vezes a partir de $N=512$, considerando só o
kernel), provavelmente porque as caches L1 e L2 das GPUs atuais já aproveitam
parte do reuso de dados da versão ingênua. Ela também fica distante da cuBLAS,
sendo cerca de 3,9 vezes mais lenta de ponta a ponta em $N=\nMax$.
Isso mostra que a vantagem da GPU depende também de otimizações difíceis de
reproduzir à mão, como as descritas por Volkov e Demmel~\cite{volkov2008}.
Da mesma forma, a nossa versão OpenMP blocada fica atrás da OpenBLAS,
cerca de 9,7 vezes mais lenta em $N=\nMax$, o que mostra o quanto a linha de base de CPU pode
ser subestimada quando não se usa uma biblioteca otimizada.}

% ===========================================================================
\section{Ameaças à Validade}\label{sec:ameacas}
% ===========================================================================

\textbf{Validade externa.} Os experimentos usam uma única combinação de CPU e
GPU; outras máquinas terão pontos de equilíbrio diferentes, embora a
Equação~\ref{eq:nstar} permita estimá-los a partir de $P_c$, $P_g$ e $B$.
Avaliamos apenas matrizes quadradas, em FP32 e com tamanhos potência de 2; GPUs
de consumo têm vazão muito menor em FP64, o que deslocaria $N^*$ para cima.
Também não medimos consumo de energia.
\textbf{Validade interna.} Frequência dinâmica (\emph{turbo}) e
estrangulamento térmico podem variar entre execuções; mitigamos com
aquecimento, repetições e o plano de energia de alto desempenho do
Windows. Os experimentos rodaram em um computador de uso pessoal, no Ubuntu sob
WSL2, com o Windows como sistema hospedeiro: processos do hospedeiro causaram
picos ocasionais de tempo, tratados com o uso da mediana, e a camada de
virtualização pode acrescentar sobrecarga às transferências entre CPU e GPU. As bibliotecas escolhem internamente seus algoritmos, o
que pode causar descontinuidades entre tamanhos.
\textbf{Validade de construto.} O tempo ponta a ponta exclui alocação e criação
de contexto; num programa que faz uma única multiplicação, esses custos
favoreceriam ainda mais a CPU. Por outro lado, aplicações que mantêm os dados
na GPU entre várias operações amortizam a transferência e se aproximam da
curva ``só kernel''.

% ===========================================================================
\section{Conclusão}\label{sec:conclusao}
% ===========================================================================

Este trabalho investigou quando a GPU compensa na multiplicação de matrizes
densas, comparando-a com uma CPU multithread otimizada e incluindo a
transferência de dados na medição. \revisar{Os resultados mostram que a GPU
só vence de ponta a ponta a partir de um ponto de equilíbrio situado
\breakeven, que o ganho chega a
\speedupEteMax$\times$ em $N=\nMax$ e que esse valor é muito menor do que o
obtido ao comparar apenas o kernel com código sequencial ingênuo.} A
mensagem prática é que a decisão de usar GPU deve considerar o tamanho do
problema, onde os dados residem e a qualidade do código de CPU disponível; o
modelo da Equação~\ref{eq:nstar} oferece uma estimativa rápida do ponto de
equilíbrio a partir de três grandezas medidas.

Como trabalhos futuros, pretendemos: avaliar cenários em que várias
multiplicações reutilizam os dados na GPU; sobrepor transferência e computação
com \emph{streams}; incluir FP64, FP16 e núcleos tensoriais; repetir os
experimentos em outras combinações de hardware; e medir a eficiência
energética.

\textbf{Disponibilidade.} \revisar{O código dos experimentos, os dados brutos e
os scripts de análise estão disponíveis em \url{https://github.com/Mth71/artigo-cpu-gpu}.}

\bibliographystyle{sbc}
\bibliography{referencias}

\end{document}
